\cotefile{91}
\scanpage{47}
pour que $\underline F_1$ ait les \unsure{vertus|conditions} dites. \deleted{\unsure{Il faut}}

$1^\circ)$ $\underline F_1$ doit être stable par multiplication par $\underline B$,
i.e. \deleted{\unsure{les|Soient}} les composantes de $b_i c_\lambda$ suivant
les $d_{\mu}$, soient $\varphi_{i\lambda\mu}\in A_{\unsure{\Phi|\varphi}}$, doivent être
\uline{nulles}, \deleted{\unsure{Il}} d'où des \uline{équations}
\[
\varphi_{i\lambda\mu}=0\qquad
\left[\begin{aligned}
1&\leq i\leq \unsure{\nu|v}\\
1&\leq\lambda\leq(n-k)\nu\\
1&\leq\mu\leq k\nu
\end{aligned}\right]
\]
% slanted
\margin{Soit $G_1$ le sous schéma de $G$ défini par \unsure{ces|ses} équations. On est ramené \unsure{à un cas}}

$2^\circ)$ Le B-module quotient \deleted{$\unsure{B|F}/\unsure{B_1|F_1}$} $F/F_1=M$
est projectif de rang \unsure{en|$n$|$k$}. Il a une
base sur $A$ de la forme $(m_i)_{1\leq i\leq k\nu}$,
Soient $\xi=(\xi_1\mathbin{-}\xi_k)$ \deleted{des} \inserted{un \unsure{syst.} de $k$} \unsure{élts|éléments} de $M$. Pour qu'ils
forment une base sur $B$, il f. et s. que
$B^k\xrightarrow{u_\xi}M$ défini soit bijectif. Or dans $B^k$ et $M$
on a des bases, donc on peut former
$\det(u_\xi)$, la condition est que $\det(u_\xi)$ soit
inversible. \deleted{\unsure{Donc pour un|Donc pour que}} Pour que
en un $\mathrm{p}^{\mathrm t}$ $x\in\operatorname{Spec}(A)$, $M_x$ soit libre sur $B_x$,
il faut et il suffit qu'il existe
$\xi\in(M_x)^k$ tel que $\det(u_\xi)$ soit inversible
dans $\deleted{\unsure{B}}A_x$. On voit facilement que
% slanted
\margin{pour éviter un \unsure{$\forall$|$\mathrm V^{\mathrm t}$}, on \unsure{pourrait expliciter des} conditions \unsure{nécessaires et suffisantes en un}}
l'on peut aussi prendre \unsure{les} $\xi\in M^k$. Donc
pour \deleted{\unsure{tel|tout}} \deleted{$\unsure{\xi\in M}$} $\xi\in M^k$, soit $U_\xi$ l'ouvert
\deleted{\unsure{en}} des non-zéros de $\det(u_\xi)$. Soit $U=\bigcup U_\xi$. C'est
% archive sheet number: 46
\scanpage{48}
l'ouvert cherché [N.B. la locale liberté de $M$ en
$x$ est en effet une condition \uline{ponctuelle}
et non seulement locale, \deleted{\unsure{c'est|exacte}} compatible avec
l'extension des corps de base $k(x)$ \ldots]

\noindent\rule{3cm}{0.4pt}\\[-0.8ex]
\rule{2.8cm}{0.4pt}

Sous les réserves habituelles, \uline{si} $g:X_{/T}\unsure{\mathbin{-}|\to}Y_{/T}$ \uline{est de
type fini}, \uline{alors} $f_\ast(g)$ \uline{est de type fini}.
[ $S$ \deleted{\unsure{est}} \uline{loc. noeth.} )

On est ramené au cas où $Y=T$. \deleted{$X$ affine
sur $T$. On peut \unsure{se supposer}} \inserted{Par}
hypothèse, \deleted{\unsure{toute|$X$}} multisection d'un
$X_s$ sur un $T_s$ est contenue dans un ouvert \inserted{affine $U$}\deleted{\unsure{aff.|$X$}}.
\unsure{Prouvons|Prenons}, qu'\unsure{une st|une qté|un nb} finie de tels ouverts
affines \deleted{\unsure{$X$}} suffisent. \deleted{\unsure{On voit que ça|Car}
suffira, dans le \unsure{cas}} \inserted{Soit en effet}
% slanted
\margin{\uline{dégager un lemme}}
\[
\deleted{\unsure{X}}\unsure{X^{\nu}|X^{v}}=
\underbrace{X\times_S\mathbin{-}\times_S X}_{\nu}\,,
\]
pour tout ouvert affine
$U$ de $X$, soit $U^\nu=\underbrace{U\times_S\mathbin{-}\times_S U}$, $U^\nu$ est un
ouvert de $X^\nu$. \deleted{\unsure{On verra que|Par}} \inserted{Comme $X^\nu$ est}
noethérien, il existe une \uline{st finie} d'ouverts
affines $U_i$ tels que $\bigcup U_i^\nu=\bigcup U^\nu$.
Soient $(x_1,\mathbin{-},x_\nu)$\deleted{$\in\unsure{\lvert X\rvert^\nu}$ \unsure{que} \unsure{les éléments d'un même ouvert affine}
$U$. Alors} soit $x\in X^{\deleted{\unsure{n}}\nu}$ de projections les $x_i$,
Alors les $x_i$ sont dans un ouvert affine
% archive sheet number: 47
\scanpage{49}
$U$ sss $x\in U^{\unsure{\nu|v}}$, \deleted{\uline{donc}} cela implique
$x\in U_i^\nu$ pour $i$ convenable i.e. les $x_\alpha\in U_i$.
O.K.

On est ainsi ramené au cas où $X$
est \uline{affine} sur $Y$. Mais alors on est
\unsure{=} ramené au cas où $X=\operatorname{Spec}(\deleted{\unsure{S}}\unsure{\mathrm V F|F|\underline F})$, où
$F$ est un faisceau loc. libre de rang \unsure{$n$}
sur $T$. On sait qu'alors
\[
f_\ast(X/T)=\operatorname{Spec}(f_\ast \unsure{\mathrm V F|F|\underline F}).
\]

\deleted{\uline{\unsure{Corollaire|Contre-exemple}} $\unsure{f_\ast}$ \unsure{Si}}

\uline{Remarque} Il n'est pas vrai en général que
la dim. des fibres de $f_{\deleted{\unsure{\ast}}\ast}(X/T)/S$
soit majorée par $\nu$ fois la dim
des fibres de \deleted{$\unsure{f_\ast}$} $X/T$. Ex:
$S=\operatorname{Spec}(k)$, $T$\deleted{$=\operatorname{Spec}\unsure{k[t]/t^2|k[t]/t^p|k[E]/t^2|k[E]/t^p}$} est purement
radiciel sur $S$, $X=T\times_S T$ i.e. $f_\ast(X/T)=$
$\underline{\operatorname{Hom}}_S(T,T)$. Ce dernier peut être de
\uline{dim} \inserted{$>0$ \deleted{$\unsure{0}$}} sur $k$. Exemple: le groupe des
automorphismes de la $k$-algèbre \deleted{$\unsure{A|L}$} est
de dim $>0$.
% slanted
\margin{\unsure{voir pour ceci|voir précisé} plus bas}
% archive sheet number: 48
\scanpage{50}
\diagram{\begin{tikzcd}[ampersand replacement=\&,column sep=large]
\& f_\ast(X/T)\arrow[r] \& f_\ast(Y/T)\\
(A) \& S'_1\arrow[u]\arrow[r,hook] \& S'\arrow[u,very thick]\arrow[ul,dashed]
\end{tikzcd}}

\diagram{\begin{tikzcd}[ampersand replacement=\&,row sep=small,column sep=small]
\& T\arrow[dl,{Straight Barb[length=1.3pt,width=2pt]}-{Stealth},thick] \& \&Y\&\&\\
S \& \&T'\arrow[ul]\arrow[dl]\arrow[ur]\arrow[rr,dashed]\&\& X\arrow[ul]\&(B)\\
\&S'\arrow[ul]\&\& T'_1\arrow[ul,hook]\arrow[dl]\arrow[ur]\&\&\\
\&\&S'_1\arrow[ul]\&\&\&
\end{tikzcd}}
\[
 T'_1\qquad\qquad T'
\]

\deleted{\uline{Proposition}} \{\uuline{Théorème?}\}
$T$ \deleted{fini et plat sur $S$ loc. noeth.} \inserted{sur $S$ (\unsure{fin.|hyp.|ou} loc. \unsure{frei|fin.|fini} \ldots) \&}
$X_{/T}\xrightarrow{g}Y_{/T}$ un morphisme \deleted{\uline{simple}} \inserted{de type fini} de T-préschémas.

On suppose que $f_\ast(X/T)$ et $f_\ast(Y/T)$ existent et
sont \deleted{loc. noethériens, que $f_\ast(X/T)\mathbin{\unsure{-|\to}}f_\ast(Y/T)$
est de type fini, que} \inserted{loc. de type fini sur $S$, et que} $S$ et $T$ sont loc. noeth.

(i) \inserted{supposons $T/S$ \unsure{affine}.} \deleted{Il faut} Si $X/T\overset{g}{\unsure{\mathbin{-}|\to}}Y/T$ est simple, il en est de $=$ de $f_\ast(g)$,

(ii) supposons $T/S$ \deleted{\unsure{fini}} \inserted{de type fini}, si $g$ est \uline{étale}, $f_\ast(g)$ est étale.

Démonstration. (i) \inserted{On peut supposer $S$ affine} Il faut pouvoir compléter le
diagramme (A) par une flèche oblique,
\deleted{\unsure{Soit|où} $S$} quand $S'_{\deleted{\unsure{1}}}$ est affine, \inserted{et de type fini sur $S$} $S'_1$ un
ss-schéma fermé ayant $=$ schéma réduit.

Cela revient à compléter de façon analogue
le diagramme (B), C'est possible, car
$T$ \unsure{étant} affine, $T'$ le sera aussi, et
$T'_1\unsure{\mathbin{-}|\to}T'$ \deleted{\unsure{une}} est une immersion fermée
définissant un isom. de $T'_1$ sur un
ss-schéma fermé de $T'$ ayant $=$ ss-schéma
\unsure{réduit que ce dernier}.
% archive sheet number: 49
\scanpage{51}
(ii) Ce raisonnement vaut aussi sans supposer $T/S$
affine, si $g$ est \uline{étale}, \deleted{\unsure{$f_\ast$}} Ainsi, $f_\ast(g)$ sera
\deleted{\unsure{ét.}}simple. Pour prouver qu'il est étale, il suffit
de prouver que \unsure{les|ses} \uline{fibres sont} \deleted{\unsure{de}} \inserted{\uline{finies}
(\unsure{les fibres|ensuite})} \deleted{de dim $0$ (\unsure{à fibres} \illegible{1})}. On est ramené
au cas où $S=\operatorname{Spec}(k)$, $k$ un corps, et
% slanted
\margin{\unsure{Donner un|Donne un|Dans un} \uline{corollaire} en fait}
\diagram{\begin{tikzcd}[ampersand replacement=\&]
\& X\arrow[r]\&Y\\
k\&T\&
\end{tikzcd}}
\diagram{\begin{tikzcd}[ampersand replacement=\&]
\& f_\ast(X)\arrow[d,very thick]\\
\operatorname{Spec}(k')\arrow[r]\& f_\ast(Y)
\end{tikzcd}}
à prouver que si $k'$ est une
clôture algébrique de $k$, alors il
n'y a qu'un nb fini de $k'$-pts de $f_\ast(X)$ au dessus d'un
$k'$-point de $f_\ast(Y)$, i.e. \deleted{\unsure{que}} il n'y a qu'un nb fini
de sections de $X'/T'$ qui
relèvent une section donnée de $Y'/T'$.
[les $'$ \unsure{dénotent|dit} ext. des scalaires de $k\unsure{:|\text{ à }}k'$].
Or $T'$ sera noethérien [de type fini sur $k'$]
donc n'a qu'un nb fini de comp. connexes.
On peut supposer $T'$ connexe. Alors le
relèvement est connu quand on
connaît l'image d'un pt. On trouve
donc

\uline{Corollaire} $1$ Soit $\unsure{\nu|v}$ le maximum des nbs
\deleted{\unsure{comp.|rangs}} de composantes connexes absolues
des $T_s$ ($s\in S$), et soit \deleted{\unsure{de}} $=$ $n$ le
maximum analogue pour $X/Y$. Alors le maximum
% archive sheet number: 50
\scanpage{52}
analogue pour $f_\ast(\unsure{X|A}/T)\unsure{\mathbin{-}|\to}f_\ast(Y/T)$ est $\leq n^{\unsure{\nu|v}}$.

\deleted{\uline{\unsure{Corol 2|Corollaire}}}

\deleted{\inserted{[Si $T/S$ est \deleted{affine} fini]} \uline{Corollaire 1} Soit \unsure{$d$ la dimension des} fibres
\unsure{de} soit $d=\sup_{y\in Y}\dim g^{-1}(y)$, et
\[
\unsure{d|d'}=\sup_{\bar y\in f_\ast(\unsure{X|Y}/T)}\dim f_\ast(g)^{-1}(\bar y).
\] \inserted{Soit $n=\deg(T/S)$,} Alors on a $d'\unsure{\leq|=}nd$,
On a égalité si \unsure{Si ttes les|toutes les} fibres
On a égalité si $\unsure{\varnothing|\delta}\deg_{\mathrm{sép}}(T/S,s)$ est indépendant
de $S$ (alors égal à $\unsure{\nu|n}$). Si les
fibres de $f$ ont ttes la $=$ dimension $d$
en tous leurs pts, il en est de $=$
\unsure{des}}

\uline{Corollaire 2} Supposons $T/S$ fini de degré $n$ partout.
Supposons qu'on est sous les conditions (i).
Soit \deleted{\unsure{de la cq}} $s\in S$, soit $D(s)$ l'ens des
dimensions des \deleted{fibres} composantes des
fibres de $X_s\unsure{\mathbin{-}|\to}Y_s$; \deleted{\unsure{de}} soit $D'(s)$ l'ens
analogue pour $f_\ast(X/T)_s\unsure{\mathbin{-}|\to}f_\ast(Y/T)_s$.
\diagram{\begin{tikzcd}[ampersand replacement=\&]
\& X_s\arrow[r]\&Y_s\\
k(s)\&T_s\&
\end{tikzcd}}
Alors $D(s)=n\,D'(s)$ En particulier,
si ttes les composantes de ttes les
fibres de $g$ ont $=$ dimension $\unsure{d}$, \inserted{alors} \deleted{\unsure{tes|toutes}}
les composantes de ttes les fibres
de $f_\ast(g)$ ont $=$ dimension $nd$.
% archive sheet number: 51
\scanpage{53}
\uline{Dém} On peut supposer $S=\operatorname{Spec}(k)$, De plus,
on peut supposer $k$ alg. clos. Enfin, d'après
\[
(\coprod f_i)_\ast(X/T)=\prod_i f_{i\ast}(X_i/T_i),
\]
on est ramené au cas où $T$ est connexe, i.e.
$T\to S$ géométriquement injectif. On peut
alors supposer $X$ et $Y$ affines, et $X\xrightarrow{g}Y$
\deleted{\unsure{un|se|en}} factorisant en $X\to\deleted{\unsure{X}}\xrightarrow{\inserted{g_1}}Y[t_1,\mathbin{-},t_d]\xrightarrow{g_2}\unsure{Y|Y'}$
\deleted{$Y\mathbin{-}X$} où $g_1$ est étale et $g_2$ la
projection can. \deleted{D'où $f_\ast(g)=f_\ast(g_2)f_\ast(g_1)$,
et $f_\ast(g_2)$ étant étale, et} Il faut prouver que
ttes les fibres de $f_\ast(g)$ ont dim $nd$.
Comme $f_\ast(g)\unsure{:|=}f_\ast(g_2)\,f_\ast(g_1)$, \inserted{et $f_\ast(g_{\unsure{1|2}})$ est étale,} on est ramené
au cas où $X=Y[t_1,\mathbin{-}t_d],=Y\times_T\unsure{\overline Z|\overrightarrow Z|\overline T}[t_1,\mathbin{-},t_d]$
\deleted{\unsure{où|i.e.} $Z=T[t$}

On est ramené (\unsure{form. de produit|form. du produit}): prouver que les fibres de
$f_\ast(T[t_1,\mathbin{-},t_d]/T)$ sont de dim $nd$,
ce qui est trivial \ldots
% archive sheet number: 52
\scanpage{54}
\diagram{\begin{tikzcd}[ampersand replacement=\&]
\&X_s\arrow[r]\&Y_s\\
k(s)=k\&T_s\&
\end{tikzcd}}
\diagram{\begin{tikzcd}[ampersand replacement=\&]
\&X'_s\arrow[r,dash]\&Y'_s\\
k'\&T'_s\arrow[u]\arrow[ur]\&
\end{tikzcd}}
\margin{Nb de comp. connexes absolues de $T_s$: $\unsure{\nu|v}_s$}
\[
\left\langle\begin{aligned}
\nu&=\sup\nu_s\\
&\hspace{-1em}\left\{\begin{aligned}n&=\deg T/S\\d&=\dim.\ \text{fibres de }g\end{aligned}\right.\\
\lambda&=\deg X/Y
\end{aligned}\right.
\]
\diagram{\begin{tikzcd}[ampersand replacement=\&]
\&\deleted{\unsure{X'_s|Y'_s}}\\
\deleted{\unsure{k'}}\arrow[r,dash]\&\deleted{\unsure{T'_s}}\arrow[u]
\end{tikzcd}}
\uline{Corollaire 1 bis} Sous les conditions du corollaire 1,
\deleted{Supposons $T/S$ fini,} \deleted{Alors} et supposons que
\deleted{\unsure{$f$ soit}} \inserted{le morphisme étale} $g:X\to Y$ soit propre de degré \deleted{$\unsure{n|d}$} $\lambda$.
Alors $f_\ast(g)_s:f_\ast(X/T)_s\unsure{\mathbin{-}|\to}f_\ast(Y/T)_s$ est
propre et de degré $\lambda^{\deleted{\unsure{n}}\nu_s}$ ($\nu_s=\deg_{\mathrm{sep}}T_s/k(s)$).

On reprend le raisonnement du cor. 1,
[bloquer cor. 3 et cor. 1 en un
seul énoncé] Il suffit de prouver
qu'il y a \uline{\unsure{exactmt|exact}} $\lambda^{\nu_s}$ points dans
l'image inverse d'un pt de $f_\ast(Y/T)$.
% margin note written sideways
\margin{Sous ces conditions, un morphisme fini étale de degré $\lambda$ est transformé \unsure{en un|en un tel|en} de degré $\lambda^\nu$}

\uline{Corollaire 3} (\deleted{$T/S$ fini} \inserted{Sous les cond. \deleted{\unsure{rappellées sur la première}} du cor 1}).
Conditions équivalentes:
\diagram{\begin{tikzcd}[ampersand replacement=\&,row sep=small,column sep=small]
\text{(i)}\arrow[d]\& \nu(s)\ \text{est \inserted{loc.} constant}\\
\text{(ii)}\arrow[d]\& \begin{gathered} X\mapsto f_\ast(X/T)\ \text{transforme morphismes \deleted{finis} étales}\\ \text{\uline{finis} en morphismes étales \uline{finis}}\end{gathered}\\
\text{(iii)}\arrow[d]\& \begin{gathered}\text{Si }T'\text{ est étale fini sur }T,\ \text{alors }f_\ast(T'/T)\\ \text{est étale fini sur }S\end{gathered}\\
\text{(iv)}\arrow[uuu,bend left=24]\& f_\ast(T\amalg T)\inserted{=\underline{\operatorname{Hom}}_S(T,S\amalg S)}\quad
\text{est étale fini sur }S\unsure{\rho|\circ}
\end{tikzcd}}
% archive sheet number: 53
\scanpage{55}
On dit alors que le morphisme \deleted{\unsure{fini}} plat
$T\unsure{\mathbin{-}|\to} S$ est \uline{homogène} \inserted{si $\unsure{\nu|v}(s)$ est loc. constant}. Cette notion est [dans le
cas de morphismes \uline{finis} du moins]
invariante par composition des morphismes,
extension de la base [trivial sur (i), qui
sera la définition]

\deleted{Corollaire}

\deleted{\uline{Proposition} Soient \deleted{$\unsure{X,Y,S}$} $T,S,X,Y$ comme
dans \unsure{ds} de Prop. précédente. Supposons
que $X$ soit un schéma en groupes sur $Y$,
et $P$ un fibré principal sur $X$}

\uline{Cor. 4} Sous les conditions de (i), supposons
$T/S$ fini, \deleted{et que $f$ \unsure{soit} \illegible{1} \unsure{l'élément}} \inserted{Si $g$ est surjectif,}
Alors $f_\ast(g)$ est \uline{surjectif}

\uline{Proposition} Le foncteur $f_\ast$ transforme
fibrés \inserted{\unsure{formellement|formellt|formelles}} principaux sur un $Y/T$
[sous réserve d'existence \inserted{et sans autre restriction}] en
fibrés \unsure{formellement|formellt|formelles} principaux.
% slanted
\margin{Pour éviter mention \unsure{compliquée} \unsure{trivialité des solutions|triviale des solutions|principal homogène}}
Il transforme fibrés principaux
simples en fibrés principaux
simples, sans supposer que $T/S$ soit fini.
Il transforme \deleted{\unsure{fibrés}} fibrés principaux à groupes
\unsure{et} \deleted{\unsure{finis|fermés}} simples \unsure{pour séparables|sont représentables|séparables} en \deleted{\unsure{des}} des
% archive sheet number: 54
\scanpage{56}
fibrés de $=$ espèce sur $S$, sous réserve que
$T\to S$ soit homogène.

En particulier, si $G$ est un groupe fini
ordinaire, un \deleted{\unsure{revêtem.}} revêtement principal de
groupe $G$ de $Y/T$ est transformé en
un \deleted{\unsure{revêtement}} fibré principal \unsure{hom|hom.|homogène} sur
$f_\ast(Y/T)$, de groupe $f_\ast(Y/T)\times_S f_\ast(\deleted{\unsure{G\times S}}T\times G/T$.
% slanted
\margin{\unsure{Pas|Pour} la peine de supposer $T/S$ fini, \unsure{pourvu que l'on sache|faisons que l'on sache} que la \unsure{descente} est possible.}

Cela montre l'intérêt de déterminer, pour
un ens fini $I$ \deleted{\unsure{variable|ordinaire}}, la variété
\deleted{\unsure{un}} étale $f_\ast(T\times I/T)=\varphi(I)$\deleted{\unsure{C'est|qui est telle que}}

\deleted{\unsure{$\varphi$} On voit que $\varphi(T)(I)$ est un foncteur
\deleted{\unsure{fini}} \uline{multiplicatif} \inserted{et $=$ exact à gauche} qui va des ens.
finis aux variétés étales de $S$. Choisissons
un pt géométrique $a$ sur $S$ supposé
connexe, on obtient ainsi un
foncteur $I\mapsto\varphi_a(T)\deleted{(I)}$ (Ens finis) dans les
ens \inserted{finis} où $\pi_1(S,a)$ opère. En tant que
foncteur à valeurs dans les ens finis,
il est de la forme $\operatorname{Hom}_{\mathrm{Ens}\,f}(A_a(T),I)$,
[petit raisonnement], et comme
$\pi_1(S,a)$ opère dans ce foncteur, il opère
dans $A_a(T)$, \deleted{qui} on a ainsi défini}
% slanted
\margin{\unsure{prendre pour principe un|pour plus de simplicité}}
% archive sheet number: 55
\scanpage{57}
\deleted{de façon naturelle un revêtement étale
\deleted{\unsure{fini de degré|$T_{\mathrm{sép}}$} \unsure{sép.} de $S$, appelé
$\underline A(T)$ ou} $\overline T=f_\ast^{\mathrm{ét}}(T)$, qui dépend foncto\uline{riellement}
de $T$, \deleted{et transforme} le foncteur $f_\ast^{\mathrm{ét}}$ transforme
d'ailleurs sommes en \deleted{produits} sommes.

On va définir un $S$-morphisme
\[
T\longrightarrow\overline T
\]
i.e. un \unsure{élt|élément} de $\operatorname{Hom}_S(T,\overline T)$, i.e. une
\uline{section} de $\underline{\operatorname{Hom}}_S(T,T$}

Plus \unsure{généralement|précisément}, pour $T$ donné, $S'\leadsto
\underline{\operatorname{Hom}}_S(T,S')$ est un foncteur des variétés
étales de $S$ dans les \unsure{var.|rev.} étales de $S$,
\uline{exact à gauche}. \deleted{Si on choisit un
pt géométrique $a$ de $S$, supposé connexe,}
\inserted{Supposons $S$ connexe.}
\deleted{On obtient ainsi un foncteur
Un théorème général nous
prouve que le foncteur est de la
forme $\varinjlim\operatorname{Hom}($}
% slanted
\margin{(\unsure{inutile?|i.e. la fonctorialité}) \unsure{raisonner|raisonnement|revient à un élément} directement sur le foncteur $S'\leadsto\operatorname{Hom}_S(T,S')$]}
Fixons un
$\mathrm p^{\mathrm t}$ \uline{géométrique} $a$ de $S$ supposé
connexe (ce ne devrait pas
être nécessaire) \deleted{\unsure{Un th. gé|Alors le foncteur}} On obtient
% archive sheet number: 56
\scanpage{58}
\margin{\unsure{!}}
Un foncteur $E\leadsto\varphi_{T,a}(E)$ des $\pi_1(S,a)$-ensembles
finis dans lui-même, foncteur qui est
exact à gauche. Je dis que un tel
foncteur est nécessairement de la
forme $\underline{\operatorname{Hom}}_{\pi_1}(A,E)$, où $A$ est
un ens, dépendant \deleted{\unsure{tel quel|bien déf.}} fonctoriellement du
foncteur $\varphi_{T,\unsure{a|c}}$ qui lui donne naissance
[produit transformé en sommes].
% slanted
\margin{à détailler, et essayer de coiffer par un th. plus général. Il y a un \unsure{pble} \unsure{ou une proposition|une variante} à établir, qui résultera de la formule analogue à $(*)$ ci-dessous, où entre une $\varinjlim$.}
On \unsure{pose=|posera|pose} $A=A(T)$, \deleted{\unsure{la|On}} et on dénote
par $f_{\unsure{\ast|!}}^{\deleted{\unsure{\ast}}\mathrm{ét}}(T)$ la variété \deleted{\unsure{séparable|équivalent}} \inserted{étale} de $S$
associée. On a donc un isom de bifoncteurs
\[
\underline{\operatorname{Hom}}_S(f_\ast^{\mathrm{ét}}(T),S')=\deleted{\unsure{H}}\underline{\operatorname{Hom}}_S(T,S')
\tag*{\makebox[0pt][r]{\makebox[\linewidth][l]{$(*)$}}}
\]
[sur $\operatorname{Rev}\,\text{ét}\,(T)\times\operatorname{Rev}\,\text{ét}\,(S)$]

\deleted{d'autres termes pour \unsure{tt|tout} $T$, le foncteur
en $S'$ du deuxième mbr est
représentable}. \unsure{Prenant|Prenons} \deleted{\unsure{les sections|donc}} \inserted{$S'=T$}, \deleted{et
\unsure{faisant passer|Prenons}} les sections des deux
mbrs, on trouve
\[
\operatorname{Hom}_S(f_\ast^{\mathrm{ét}}(T),S')\simeq\operatorname{Hom}_S(T,S')
\]
d'où, en faisant $S'=f_\ast^{\unsure{\mathrm{ét}}}(T)$, \unsure{on trouve|un hom}
% archive sheet number: 57
